\chapter{\texttt{pipeline.solar\_earth} --- solar production composed with Earth regeneration}
\label{ch:pipeline-solar-earth}

\texttt{pipeline.solar\_earth} composes \cref{ch:pipeline-solar,ch:pipeline-earth}
into the single workflow most solar-neutrino analyses actually need:
production inside the Sun, incoherent propagation to the Earth (the mass
mixture decoheres over the Sun--Earth distance, \cref{ch:core-common}),
then Earth-crossing regeneration to a detector. It introduces exactly one
piece of genuine composition logic beyond simply calling the two pipelines
in sequence: combining the Sun--Earth distance modulation with the
detector-exposure choice, derived below.

\begin{theorybox}
\subsection*{Combining the Sun--Earth distance modulation with detector exposure}
Solar-model flux tables (\cref{ch:solar}) are normalised to 1~AU, but
Earth's elliptical orbit makes the physically received flux vary by about
$\pm3.4\%$ over the year, via the instantaneous factor
$(1\,\mathrm{AU}/R(\text{date}))^2$
(\pyfunc{medium.solar.geometry.sun\_earth\_distance\_factor}). Exactly one
of two ways to apply it is meaningful for a given call, matching whether
the detector-exposure average (\cref{sec:earth-exposure}) is active:
\begin{itemize}
  \item \textbf{No exposure averaging} (\code{integrate\_exposure} resolves
    to \code{False}): pass an explicit ISO \code{date="YYYY-MM-DD"}: the
    instantaneous factor for that single calendar date is applied to
    \code{detector\_flux}. Passing \code{date} together with an
    exposure-averaged call raises, since a single date is not meaningful
    once already averaging over a day-of-year window.
  \item \textbf{Exposure averaging active} (\code{integrate\_exposure}
    resolves to \code{True}): pass \code{average\_sun\_earth\_distance=True}
    instead, which averages the factor uniformly over the \emph{same}
    day-of-year window already used for the nadir-angle exposure average
    (\code{config.exposure.exposure\_d1}/\code{exposure\_d2}, via
    \pyfunc{sun\_earth\_distance\_factor\_averaged}), so one exposure window
    consistently accounts for both the detector's day/night geometry and
    the Sun--Earth distance over the same period. Passing it without
    exposure averaging active raises, for the symmetric reason: there is no
    day-of-year window to average over.
\end{itemize}
Omitting both (the default) leaves the result at the solar model's native
1~AU normalisation, exactly as \cref{ch:pipeline-solar} alone would.
\end{theorybox}

\section{\texttt{pipeline.solar\_earth} --- composed workflow}

\modulehead{pipeline/solar\_earth.py}{\pyfunc{propagate\_solar\_to\_earth\_detector}:
solar production and incoherent Earth regeneration, without duplicating the
intermediate mass-mixture state.}

\begin{implbox}
\begin{itemize}
  \item \pyclass{SolarEarthDetectorResult} (frozen dataclass): the full
    \code{solar\_surface} (\pyclass{SolarSurfaceResult},
    \cref{ch:pipeline-solar}) and \code{earth\_detector}
    (\pyclass{EarthDetectorResult}, \cref{ch:pipeline-earth}) sub-results,
    plus the composed \code{detector\_flux}, \code{detector\_probabilities}
    (an alias into whichever of \code{earth\_detector.probabilities\_eta}/%
    \code{probabilities\_exposure} was populated -- \textbf{the same tensor
    object}, not a copy), and the optional energy-integrated
    \code{probabilities\_energy\_averaged}/%
    \code{detector\_flux\_energy\_integrated}.
  \item \pyfunc{propagate\_solar\_to\_earth\_detector(*, E\_MeV, config,
    source, solar\_profile=None, earth\_profile=None, eta=None,
    source\_spectrum=None, integrate\_exposure=None, integrate\_energy=False,
    solar\_method="adiabatic\_approximated", legacy\_precision=False,
    include\_matter\_nc=None, date=None, average\_sun\_earth\_distance=False)}:
    calls \pyfunc{propagate\_solar\_to\_surface} (forwarding
    \code{solar\_method}/\code{include\_matter\_nc} for the solar leg only
    -- the Earth leg has no analogous parameter exposed here yet), feeds
    \code{solar.mass\_weights} \emph{directly} into
    \pyfunc{propagate\_earth\_to\_detector}
    (\code{incident\_basis="mass"}) or its \code{\_exposure} variant
    according to \code{integrate\_exposure} (\code{None} resolves from
    \code{config.exposure.integrate\_exposure}), then applies
    \pyfunc{core.common.flux.flux\_state} with the source's total flux and
    spectrum (\code{source\_spectrum}, or \code{profile.spectrum(source,
    \ldots)} when omitted) to obtain \code{detector\_flux}. Setting
    \code{integrate\_energy=True} additionally reduces both the
    probabilities and the flux over the energy grid via
    \pyfunc{probability\_integrated}/\pyfunc{flux\_integrated}
    (\cref{ch:core-common}).
\end{itemize}
\end{implbox}

\begin{examplebox}[title={Boron-8 solar neutrinos detected on a nadir-angle grid}]
\begin{lstlisting}[style=tpython]
import torch

from tpeanuts.config.propagation import PropagationConfig
from tpeanuts.pipeline.solar_earth import propagate_solar_to_earth_detector
from tpeanuts.util.context import RuntimeContext

context = RuntimeContext.resolve("cpu", torch.float64)
oscillation = PropagationConfig.oscillation_parameters_from_preset(
    "_SM_NUFIT52_NO", context=context,
)
config = PropagationConfig(runtime=context, oscillation=oscillation)

result = propagate_solar_to_earth_detector(
    E_MeV=[5.0, 8.0],
    config=config,
    source="8B",
    eta=[0.4, 1.0],
    integrate_exposure=False,
)

# The Earth leg's incident state IS the solar mass-mixture, not a copy:
assert result.earth_detector.incident_state is result.solar_surface.mass_weights
print(result.detector_flux.shape)   # torch.Size([2, 2, 3]): (E, eta, flavour)
\end{lstlisting}
\end{examplebox}

Requesting the exposure-averaged detector rate instead only changes two
arguments -- \code{eta} is dropped and \code{integrate\_exposure=True} takes
its place, resolving to \pyfunc{propagate\_earth\_to\_detector\_exposure}
internally -- and \code{integrate\_energy=True} additionally collapses the
energy axis:
\begin{examplebox}[title={Exposure- and energy-integrated event rate}]
\begin{lstlisting}[style=tpython]
spectrum = torch.tensor([0.4, 0.6], dtype=torch.float64)

rate_result = propagate_solar_to_earth_detector(
    E_MeV=[4.0, 8.0],
    config=config,
    source="8B",
    source_spectrum=spectrum,
    integrate_exposure=True,
    integrate_energy=True,
    average_sun_earth_distance=True,
)
print(rate_result.detector_flux_energy_integrated.shape)   # torch.Size([3])
\end{lstlisting}
\end{examplebox}
